\section{Energy, Enstrophy, And The Source-Reserve Branch}
\label{sec:energy-enstrophy-source-branch}

The first positive-forward branch starts from the strongest native control in the
Navier--Stokes equation: kinetic energy. This branch is the natural starting
point because it uses the equation on its own surface. There is no imported
geometry, no auxiliary model, and no change of solution class. The original
solution pays an exact bill:
\[
  \frac12\frac{d}{dt}\|u(t)\|_{L^2}^2
  +\nu\|\nabla u(t)\|_{L^2}^2=0.
\]
After integration, this gives
\[
  \sup_{0\le t<T}\|u(t)\|_{L^2}^2
  +2\nu\int_0^T\|\nabla u(t)\|_{L^2}^2\,dt
  \le \|u_0\|_{L^2}^2.
\]
The estimate is exact, global on every finite time interval, and same-solution.
It also shows the limitation of the branch. Smooth continuation needs control of
derivatives beyond the energy level.

The natural lift is enstrophy. Define
\[
  Y(t)=\frac12\|\nabla u(t)\|_{L^2}^2.
\]
For smooth solutions on the periodic surface, differentiating and integrating by
parts gives
\[
  \frac12\frac{d}{dt}\|\nabla u(t)\|_{L^2}^2
  +\nu\|\Delta u(t)\|_{L^2}^2
  =
  \int_{\mathbb T^3}(u\cdot\nabla)u\cdot\Delta u\,dx.
\]
The left side is the derivative-level viscous payment. The right side is the
source term. It is the place where the motion of the fluid can feed the same
derivatives that viscosity is trying to damp.

The standard estimate is useful and sharp enough to explain the obstruction.
H\"older gives
\[
  \left|\int_{\mathbb T^3}(u\cdot\nabla)u\cdot\Delta u\,dx\right|
  \le
  \|u\|_{L^6}\|\nabla u\|_{L^3}\|\Delta u\|_{L^2}.
\]
Sobolev and interpolation give
\[
  \|u\|_{L^6}\lesssim \|\nabla u\|_{L^2},
  \qquad
  \|\nabla u\|_{L^3}
  \lesssim
  \|\nabla u\|_{L^2}^{1/2}\|\Delta u\|_{L^2}^{1/2}.
\]
Therefore
\[
  |I(t)|
  \lesssim
  \|\nabla u\|_{L^2}^{3/2}\|\Delta u\|_{L^2}^{3/2}.
\]
Young's inequality yields
\[
  |I(t)|
  \le
  \frac{\nu}{2}\|\Delta u\|_{L^2}^2
  +C_\nu\|\nabla u\|_{L^2}^6,
\]
and hence
\[
  Y'(t)\le C_\nu Y(t)^3.
\]

This is a continuation estimate, but it is not a global smoothness estimate. The
comparison equation
\[
  Z'(t)=C_\nu Z(t)^3
\]
has finite-time growth. The branch therefore stops at a precise point. Energy
pays the integrated first-derivative bill. Enstrophy exposes the nonlinear source
bill. The known absorption gives local control and leaves a source channel that
can still support terminal derivative growth.

The survivor of this branch is the source reserve: the part of the nonlinear
interaction that remains after the available viscous payment has been made. In a
positive-forward proof, one tries to delete this reserve by a sharper estimate,
a cancellation, a sign argument, a pressure decomposition, a localized criterion,
or a scale separation. Each refinement may reduce the reserve. The branch remains
open whenever a terminal source object can still be carried by the same solution.

For the rebuilt paper, this is the branch's main lesson. The obstruction is not a
failure of notation. It is a terminal object candidate. A finite-time breakdown
claim that relies on this source channel has to say what the channel is at the
terminal time. If it still participates in the same Navier--Stokes equation, it
belongs to the same-solution witness test. If the source reserve breaks
same-equation participation, loses the carrier required for the selected packet,
or fails positive-scale field coherence, then the obstruction is read through the
CM-exit method rather than treated as an in-class nonsmooth continuation.

Thus the energy/enstrophy branch supplies one of the central motivating families
for the class-membership proof. The pass side is the ordinary continuation side:
the source is absorbed strongly enough to propagate the derivative control. The
fail side is a terminal source-reserve object. Its final status is determined by
the same-solution membership requirements, not by the mere fact that a positive
estimate stalled.
